\section{Introduction}
A small language model given a text editor as its only tool fails at editing tasks that a larger model solves reliably.
It chooses a plausible command, the editor reports an error, and the model sends the same command again until its
budget is gone. This paper is about taking such a model from that baseline to competence: which data, collected how,
and trained with which signal, turn a model that solves none of 100 held-out editing tasks into one that solves a
quarter to a third of them.

The task is editing files with headless Vim. A request names an outcome (``make every status line in the log use the
new field order''), the agent has one tool that runs Ex commands in a fresh Vim and returns what happened, and a
hidden check grades the files it leaves behind. The tasks are small and fully verifiable, which is what makes the
process below possible: every trajectory a model produces can be graded, so every trajectory is either a verified
solution or a verified failure.

The paper's scope is the process, from baseline to competence for models of 0.6B, 4B and 8B parameters, and its
measurement on held-out tasks. It contributes:
\begin{itemize}
\item \textbf{A data process} that starts from verified teacher trajectories, adds the student's own verified
  solutions, and adds recovery data: a teacher continues the student's failed attempts from the step where they went
  wrong, so the student sees its own error followed by a correction. New records enter the training set only as new
  solutions of their task, and every record names the model that made it.
\item \textbf{A measurement} on 100 held-out exams under a fixed protocol: one attempt per task, then up to four
  reruns of the tasks that failed, a paired test over tasks, and validity rules for infrastructure failures.
\item \textbf{A loop penalty} in a policy-gradient reward, aimed at the dominant failure of small models, repeating a
  failed call until the loop watchdog ends the run, and the finding that it does not by itself reduce looping.
\item \textbf{A reproduction recipe} that rebuilds every number from released tasks and traces with one orchestrator
  (\cref{app:repro}).
\end{itemize}
\Cref{sec:task} defines the task and the measurement, \cref{sec:method} the process, \cref{sec:results} the results
and \cref{sec:discussion} what they do and do not show.
